\section*{Anexo 1 de la Fase 1: definición formal de los formatos}
\label{sec:F1A01}

\subsection*{Nota previa}
\label{sec:F1A01Nota}

La definición de los formatos no se incluyó de forma formal y
completa durante las iteraciones anteriores: a medida que el proyecto las
necesitaba los formatos se fueron añadiendo, ampliando y afinando en cada iteración
(el setup al comenzar la Fase 1, el GWB al empezar la autoría, el GWI al aparecer el primer guardado). Este anexo se adjunta al final de la iteración como referencia:
recoge la definición rigurosa y completa de los tres formatos tal y como están
implementados, y sirve de contrato para las fases siguientes.

\subsection*{Conceptos comunes a los tres formatos}
\label{sec:F1A01Comun}

Todo documento JSON del proyecto comparte dos datos comunes, la pareja de
identificación: \texttt{tipo} (GWS para el setup, GWB para el cartucho base, GWI para el 
cartucho de autoría) y \texttt{version} con un numero que la identifica. 

El campo \texttt{tipo} del JSON distingue el
formato. Cuando la carga de un cartucho recibe una carpeta o un ZIP con varios
\texttt{.json}, se elige uno por prioridad de tipo: GWC, luego GWI, luego
GWB, avisando por consola de la elección.

La validación de entrada (el contrato de lo que puede ENTRAR) exige incluir en todos los 
formatos de cartucho estos datos básicos:

\begin{itemize}
	\item \texttt{Tipo} presente y válido (GWB, GWI o GWC): sin tipo no se
	asume nada.
	\item \texttt{Version} para poder cargar datos correctamente según la versión del fichero (actualmente todos tienen versión 1).
	\item \texttt{pantalla} con las tres imágenes: \texttt{img\_s} OBLIGATORIA (de ella se extrae la geometría) e \texttt{img\_f} e \texttt{img\_m} OPCIONALES. Si una imagen se declara, su fichero DEBE estar presente; sin fondo, el simulador pinta la pantalla en blanco. (El GWC es binario y se valida por su propio camino, en la lectura del fichero.)
\end{itemize}

\subsection*{El formato GWS: el setup de la máquina}
\label{sec:F1A01GWS}

El GWS \footnote{En la Fase 1 se usó como nombre general setup, pero en la fase 5 se decidió unificar y denominarlo GWS} (Game \& Watch Setup) es el JSON que define el HARDWARE destino final: describe la
máquina fija sobre la que se cargan y ejecutan finalmente los cartucho. 

No es un cartucho: no lleva imágenes ni juego, y se carga con su propio botón (Cargar
Setup) como un único fichero \texttt{.json}. Recargar el setup
reinicia pantalla y botones y limpia el cartucho en memoria.

Los campos del JSON del setup:

\begin{table}[H]
	\centering
	\begin{tabular}{llp{10cm}}
		\toprule
		Campo & Tipo & Contenido \\
		\midrule
		\texttt{tipo} & cadena & \texttt{"GWS"}\footnotemark \\
		\texttt{version} & cadena & Versión del formato (actualmente \texttt{"1.1"}) \\
		\texttt{nombre} & cadena & Nombre de la máquina \\
		\texttt{estilo} & cadena & Estilo de la máquina (\texttt{"2"}, \texttt{"4"} o \texttt{"5"}): elige el esqueleto físico del simulador \\
		\texttt{pantalla} & objeto & \texttt{\{ancho, alto\}} en píxeles de la pantalla LCD \\
		\texttt{botones} & array & Un \texttt{\{numero, puerto, alias?, tecla?\}} por botón físico \\
		\bottomrule
	\end{tabular}
	\caption{Campos del json del setup GWS}
	\label{tab:F1A01GWS}
\end{table}
\footnotetext{Inicialmente se definió como GW, en fase 5 se unificó a GWS.}


Los campos de cada botón:

\begin{itemize}
	\item \texttt{numero}: el número físico del botón, cadena
	\texttt{"0"}--\texttt{"8"} (se aceptan números sin comillas y se
	convierten con aviso). Los números 6, 7 y 8 son el trío de sistema
	(A / B / TIME).
	\item \texttt{puerto}: el GPIO del ESP32 al que va conectado. El
	simulador no lo usa: es el dato que leerá el firmware de la Fase 6.
	\item \texttt{alias}: OPCIONAL. Nombre corto de sistema para los
	botones 6, 7 y 8 (\texttt{A}, \texttt{B}, \texttt{R}).
	\item \texttt{tecla}: OPCIONAL. La letra del teclado del PC que
	acciona el botón; si falta, se aplica la tabla por defecto del
	motor.
\end{itemize}

\begin{lstlisting}[
	language=json, 
	frame=single, 
	breaklines=true, 
	basicstyle=\ttfamily\small,
	caption={Ejemplo de json para un setup GWS},
	label={lst:setup_gws}
	]
	{
		"tipo": "GWS",
		"version": "1.1",
		"nombre": "Parachute",
		"estilo": "2",
		"pantalla": { "ancho": 156, "alto": 156 },
		"botones": [
		{"numero": "0", "puerto": 4, "tecla": "A"},
		{"numero": "4", "puerto": 5, "tecla": "L"},
		{"numero": "6", "puerto": 18, "alias": "A", "tecla": "1"},
		{"numero": "7", "puerto": 19, "alias": "B", "tecla": "2"},
		{"numero": "8", "puerto": 21, "alias": "R", "tecla": "3"}
		]
	}
\end{lstlisting}

\subsection*{El formato GWB: la base de trabajo}
\label{sec:F1A01GWB}

El cartucho de tipo GWB (Game \& Watch Base) es el cartucho de PARTIDA de la autoría: contiene sobre todo la imagen de partida con los objetos dibujados, todavía sin
segmentos identificados. Se \emph{entrega como una CARPETA} (no un ZIP) con estos ficheros en su interior:

\begin{itemize}
	\item Un \texttt{.json} con los datos del cartucho GWB (cualquier nombre; si hay
	varios, la prioridad del apartado anterior).
	
	\item \texttt{pantalla.img\_f}: OPCIONAL. La imagen de fondo sin los segmentos. Hay 
	cartuchos sin arte de fondo. Si \texttt{img\_f} está declarada, el 	fichero referenciado
	DEBE estar presente en la carpeta; sin fondo, el simulador pinta la pantalla en blanco.
	
	\item \texttt{pantalla.img\_s}: OBLIGATORIO. La imagen con todos los segmentos encendidos pero sin 
	el fondo, OBLIGATORIA siempre (de ella se extrae la geometría): el fichero DEBE estar presente en la carpeta. Si existe \texttt{img\_f}
	debe medir exactamente lo mismo; si no es así hay que dar un aviso por consola NO bloqueante.
	
	\item \texttt{pantalla.img\_m}: OPCIONAL. La portada de la máquina, que se usará en la
	pantalla de selección de juegos como referencia gráfica del cartucho a cargar. Si se declara, el fichero referenciado DEBE estar presente.
\end{itemize}

Los campos del JSON asociado a un GWB:

\begin{table}[H]
	\centering
	\begin{tabular}{llp{10cm}}
		\toprule
		Campo & Tipo & Contenido \\
		\midrule
		\texttt{tipo} & cadena & \texttt{"GWB"} \\
		\texttt{version} & número & Versión del formato (actualmente 1) \\
		\texttt{nombre} & cadena & Nombre del cartucho \\
		\texttt{estilo} & cadena & Estilo de la máquina (debe coincidir con el setup) \\
		\texttt{pantalla} & -- & Tres campos \texttt{\{img\_m, img\_f, img\_s\}} con los  nombres de los ficheros de imagen \\
		\texttt{botones} & array & OPCIONAL: Los nombres que damos a los botones. \texttt{[\{numero, id\}]} por cada botón \\
		\bottomrule
	\end{tabular}
	\caption{Campos del json del GWB}
	\label{tab:F1A01GWB}
\end{table}

Un GWB NO lleva sprites, segmentos ni autoría: al cargarlo, el editor
descompone \texttt{img\_s} automáticamente (extracción de componentes
negras) y el diseñador trabaja a partir de ahí. 

El campo \texttt{numero} de un botón es el número físico del hardware (los
\texttt{id} son etiquetas legibles como \texttt{IZQ} o \texttt{Start}).

\begin{lstlisting}[
	language=json, 
	frame=single, 
	breaklines=true, 
	basicstyle=\ttfamily\small,
	caption={Ejemplo de json para un cartucho GWB},
	label={lst:setup_cartucho}
	]
	{
		"tipo": "GWB",
		"version": 1,
		"nombre": "Parachute",
		"estilo": "2",
		"pantalla": {
			"img_m": "parachute_m.png",
			"img_f": "parachute_f.png",
			"img_s": "parachute_s.png"
		},
		"botones": [
		{"numero": 0, "id": "IZQ"},
		{"numero": 4, "id": "DER"},
		{"numero": "A", "id": "Start"}
		]
	}
\end{lstlisting}

\subsection*{El formato GWI: el cartucho con autoría}
\label{sec:F1A01GWI}

El GWI (Game \& Watch Intermedio) es el cartucho con los datos del juego y los comandos efectuados, mas los datos de intercambio: guarda las listas que produce el editor (sprites y segmentos),
las imágenes desglosadas y el HISTORIAL de la autoría. Se guarda como un ZIP
(\texttt{cartucho\_gwi.zip}), aunque se puede descomprimir y usar la carpeta generada \footnote{Esto se decide realmente en la Fase 5 Iteración 9 (ver sección \ref{sec:F5I9Usuario})} con esta disposición (el simulador lo carga
IGUAL en ZIP que descomprimido en carpeta: la forma la detecta el
programa por el contenido, no el botón):

\begin{itemize}
	\item \texttt{cartucho\_gwi.json} en la raíz: el JSON del cartucho.
	\item La imagen de fondo, en la raíz, con el nombre \texttt{pantalla.img\_f}. OPCIONAL, igual que en el GWB: los cartuchos sin arte de fondo no la llevan y el ZIP prescinde de ella; sin fondo, el simulador pinta la pantalla en blanco.
	\item La imagen de segmentos, en la raíz, con el nombre 
	\texttt{pantalla.img\_s}. OBLIGATORIA también en el GWI (mismo contrato que el GWB): si el ZIP no la lleva, el GWI no carga; al guardar se avisa en consola si no hay imagen de segmentos en memoria.
	\item Carpeta \texttt{segmentos/} con UN PNG por cada sprite, llamado
	\texttt{img\_NNN.png} (y \texttt{img\_FNN.png} los de las piezas
	ficticias); el campo \texttt{fichero} de cada sprite referencia su
	PNG, y la resolución por nombre base hace indiferente la carpeta al
	cargar.
\end{itemize}

El json es un formato que permite ignorar estacios y saltos de línea no significativos, por lo que se genera de forma compacta para mejorar su lectura, si no el tamaño puede ser muy grande para una cómoda visualziación.

Los campos incluidos en el json son:

\begin{table}[H]
	\centering
	\begin{tabular}{llp{10cm}}
		\toprule
		Campo & Tipo & Significado \\
		\midrule
		\texttt{tipo} & cadena & \texttt{"GWI"} \\
		\texttt{version} & número & Versión del formato (1) \\
		\texttt{nombre} & cadena & Nombre del cartucho \\
		\texttt{estilo} & cadena & Estilo de la máquina \\
		\texttt{codigoGWC}\footnotemark & cadena &  OPCIONAL: El Código GWC del par, un texto con formato \texttt{"\#HHHH\#"}. Se guarda al generar el GWC y GWI a la vez, si solo se guarda el GWI se genera sin este campo\\
		\texttt{pantalla} & objeto & \texttt{\{img\_m, img\_f, img\_s\}} (igual que el GWB) \\
		\texttt{fin} & objeto & OPCIONAL: \texttt{\{mecanismo, contador\}} con \texttt{mecanismo} en \texttt{vida}, \texttt{muerte} o \texttt{contador} (contador solo con el tercero, 1--4). Propiedad del cartucho, no un segmento \\
		\texttt{secuencia} & objeto & OPCIONAL: \texttt{\{ticks, segmentos\}}: animación periódica, hasta 5 segmentos lanzados en orden, uno por tick \\
		\texttt{pausas} & objeto & OPCIONAL: \texttt{\{simulador, externo\}} en ms; ausente o sin pausa: 500 ms \\
		\texttt{aceleracion} & objeto & OPCIONAL: \texttt{\{cada, resta, minimo\}}: cada \texttt{cada} ticks la pausa vigente baja \texttt{resta} ms sin pasar de \texttt{minimo} \\
		\texttt{rutas} & objeto & OPCIONAL: \texttt{\{aumentarCada\}}: cada N rutas lanzadas el máximo de rutas simultáneas sube en 1 (ausente: máximo 1) \\
		\texttt{botones} & array & \texttt{[\{numero, id\}]} por botón \\
		\texttt{sprites} & array & Los sprites, ver más abajo \\
		\texttt{segmentos} & array & Los segmentos, ver más abajo \\
		\texttt{autoria} & objeto & El historial de la autoría, ver más abajo \\
		\bottomrule
	\end{tabular}
	\caption{Campos del GWI}
	\label{tab:F1A01GWI}
\end{table}
\footnotetext{Este campo se incluyó en la Fase 5 para tener emparejados GWI y GWC.}


Los campos opcionales usan los valores por defecto del motor cuando no
viajan: el GWI no lleva sobras (los defaults no se escriben).

\subsubsection*{Los segmentos}
\label{sec:F1A01Segmentos}

Cada segmento lleva los campos comunes \texttt{\{id, x, y, sprite, tipo,
	estado, cambiado\}}: identificador \texttt{seg\_NNN}, posición en
coordenadas de la imagen fuente, el sprite que lo dibuja, el tipo, y el
estado y el indicador de cambio que el motor de estados escribe en cada
tick (la lista ES el estado). Los enlaces válidos dependen del tipo:

\begin{table}[H]
	\centering
	\begin{tabularx}{\textwidth}{l l l X}
		\toprule
		Nº & Tipo & Enlaces propios & Otros campos \\
		\midrule
		0 & jugador & \makecell[tl]{\texttt{izquierda,}\\ \texttt{derecha,}\\ \texttt{arriba,}\\ \texttt{abajo}} & \texttt{primero} (jugador inicial),\newline \texttt{colision},\newline \texttt{salvadoContador},\newline \texttt{muerteContador} \\
		
		\cmidrule(lr){1-4}
		
		1 & juego & \makecell[tl]{\texttt{siguiente,}\\ \texttt{salvado,}\\ \texttt{muerte,}\\ \texttt{mostrarJunto}} & \texttt{ruta} (nº de ruta),\newline \texttt{contador} \\
		
		\cmidrule(lr){1-4}
		
		2 & inicio & como juego & El arranque de una ruta \\
		
		\cmidrule(lr){1-4}
		
		3 & vida & como juego & \texttt{contador} (orden en la cadena de vidas) \\
		
		\cmidrule(lr){1-4}
		
		4 & muerte & como juego & \texttt{contador} \\
		
		\cmidrule(lr){1-4}
		
		5 & dígito & \texttt{barra\_a}\ldots{}\texttt{barra\_g} & Las siete barras del «8» \\
		
		\cmidrule(lr){1-4}
		
		6 & contador & \texttt{digito\_1}\ldots{}\texttt{digito\_4} & \texttt{contadorInterno} (1--4),\newline \texttt{valorInicial},\newline \texttt{direccion} (ascendente/descendente),\newline \texttt{presentarCero} \\
		
		\cmidrule(lr){1-4}
		
		7 & cartucho & --- & La ficha de datos generales del cartucho: NO viaja en la lista (su mecanismo de fin es el campo \texttt{fin} del nivel superior) \\
		
		\cmidrule(lr){1-4}
		
		8 & adicional & \texttt{siguiente} & Pieza adicional \\
		
		\cmidrule(lr){1-4}
		
		9 & inicioAdicional & \texttt{siguiente} & Arranque de una cadena adicional \\
		\bottomrule
	\end{tabularx}
	\caption{Tipos de segmento y sus enlaces válidos}
	\label{tab:F1A01Tipos}
\end{table}

A ellos se añaden los campos transversales: \texttt{visible} (falso en los
segmentos ocultos: absorbidos y piezas internas), \texttt{fusionadoPor} (el
identificador de la pieza ficticia que lo absorbió: viaja \texttt{marcado} y el
compilador del GWC lo eliminará), y los \texttt{inversos} calculados por la
verificación al guardar: \texttt{colisionDe} (inverso de \texttt{colision})
y \texttt{mostradoJunto} (inverso de \texttt{mostrarJunto}). 

Los inversos
viajan de forma deliberadamente redundante: el compilador trabaja en el PC
sin cálculo extra. Los enlaces apuntan a identificadores legibles
(\texttt{seg\_NNN} o \texttt{FNN}); un enlace nulo (\texttt{null}) significa
sin enlace.

\subsubsection*{Los sprites}
\label{sec:F1A01Sprites}

Cada sprite es \texttt{\{id, ancho, alto, fichero\}}: el identificador
legible (\texttt{img\_NNN}, o \texttt{img\_FNN} para las piezas ficticias),
sus dimensiones en píxeles de la imagen fuente y la ruta de su PNG dentro
del ZIP. Los sprites se comparten POR IDENTIDAD cuando varios segmentos
nacen de la misma componente o pieza: la deduplicación de imágenes es parte
del formato desde el origen (una imagen, una entrada en la lista, un PNG).


\subsubsection*{La autoría: el historial}
\label{sec:F1A01Autoria}

La sección \texttt{autoria} guarda \texttt{\{cambios, ultimoOp\}}: el
registro CRONOLÓGICO de acciones es la única fuente de verdad de la
autoría. Cada entrada lleva al menos \texttt{\{n, relacionados\}} (número de
acción, que CONTINÚA entre guardados, y nº de elementos tocados) más los
campos de la acción con sus valores nuevos; el par
\texttt{\{segmento, campo, antes\}} describe el cambio y \texttt{antes}
conserva el valor previo (\texttt{null} si el campo nace). Al cargar un GWI,
el historial se REPRODUCE en orden sobre el estado de nacimiento (la
geometría extraída de \texttt{img\_s}): así el fichero es replay-safe y
deshacer/rehacer funciona tras recargar. Las altas y bajas de piezas
ficticias viajan como acciones (\texttt{altaFicticio}, \texttt{bajaFicticio})
con su caja y sus componentes absorbidas.

\subsubsection*{Serialización}
\label{sec:F1A01Serializacion}

El JSON se escribe legible pero COMPACTO: los arrays largos (sprites,
segmentos, cambios) van a UNA línea por elemento, el resto de campos a una
línea por campo, todo en UTF-8. 

El objetivo es doble: que el fichero sea
diff-friendly (una acción de la autoría, una línea del historial) y que un
diseñador pueda de manera sencilla crearlo o editarlo a mano, si lo ve 
necesario (aunque no se recomienda, el sistema está pensado para usar el 
entorno gráfico que permite definir el juego sin necesidad de editar el json).